\if1\anon
{
  \title{\bfseries Large series priors for adaptive and optimal Bayesian nonparametric regression}
  \author{Felipe Bedoya \\
    Department of Statistics, Rice University\\
    and \\
    Meng Li \\
    Department of Statistics, Rice University}
  \date{}\maketitle
} \fi

\if0\anon
{
  \bigskip
  \bigskip
  \bigskip
  \begin{center}
    {\LARGE\bfseries Large series priors for adaptive and optimal Bayesian nonparametric regression}
\end{center}
  \medskip
} \fi

\bigskip
\begin{abstract}
Random series priors are widely used in nonparametric Bayesian regression, and their performance depends critically on the number of basis functions. A fixed truncation level is simple but must be matched to the unknown smoothness, while hierarchical priors on the number of terms are adaptive but require sampling over spaces of varying dimension. We propose a large series prior: a deterministic number of basis functions growing as $n^b$ for a constant $b>0$, where $n$ is the sample size, combined with a shrinkage prior on the coefficients. We establish posterior contraction rates for product shrinkage priors satisfying prior-mass conditions, with general dictionaries; the rate is governed by how sparsely the dictionary represents the regression function. For anisotropic H\"older-type classes and Fourier or B-spline dictionaries, one prior, fixed in advance, attains near-minimax rates pointwise over the range of smoothness allowed by the dictionary; the Horseshoe with a deterministic global scale satisfies the required conditions. Posterior computation is that of a linear regression of fixed dimension under a shrinkage prior, so standard samplers apply. In simulations, enlarging the dictionary leaves the Horseshoe error essentially unchanged, unlike ridge regression, and the empirical error decay follows the sparsity of the fitted coefficients, as the theory predicts. Three real-data applications illustrate the method.
\end{abstract}

\noindent%
{\it Keywords:} Bayesian nonparametrics, random series prior, adaptivity, shrinkage priors, sparsity, posterior contraction, minimax optimality 
\vfill
